% ================================================================
%  chapters/parts/ch09-b.tex
%  第 9 章 表面分析技术基础（分片 b）
%  源：S2（下册）§9.2.5—§9.2.6，书内 p.615—623 / PDF p.106—114
%  覆盖：电子结合能、初态效应与化学位移；终态效应及伴峰
%  承接分片 a（9.1 与 9.2.1—9.2.4）；分片 d 续写 9.2.7—9.2.10
% ================================================================

\subsection{电子结合能、初态效应与化学位移}

\begin{probout}
XPS 最核心的可测量是\textbf{电子结合能} $E_{\mathrm{B}}$，它反映原子中电子与核电荷之间的
相互作用强度。同一元素在不同化学环境中的 $E_{\mathrm{B}}$ 差异称为\textbf{化学位移}，
这是 XPS 判断化学态的依据。但结合能并不是初态的性质：它取决于电离前后
\textbf{整个体系能量之差}，因此既有初态贡献也有终态贡献。
分清这两者，是判断一个位移能否被解释为「价态变化」的前提。
\end{probout}

\subsubsection{结合能的定义与光电发射的三步模型}

结合能定义为单位过程前后体系能量之差：体系初态 $i$ 有 $n$ 个核外电子，
终态 $f$ 被电离、有 $n-1$ 个电子外加一个自由电子，则

\begin{equation}\label{eq:eb-def}
E_{\mathrm{B}}=E_f(n-1)-E_i(n)
\end{equation}

式~\eqref{eq:eb-def} 说明 $E_{\mathrm{B}}$ 是\textbf{两个多电子态之间的能量差}，
而不是单个轨道的「轨道能」。因此任何改变初态或终态能量的因素都会改变测得的结合能。

光电子发射可分解为三步，每一步都对谱结构有贡献：

\begin{center}
{\srcsize
\begin{tabularx}{\columnwidth}{@{}lX@{}}
\toprule
步骤 & 对谱的贡献 \\
\midrule
吸收光子能量并电离 & \textbf{初态效应}：决定主峰位置（化学位移的主要来源） \\
原子响应与光电子产生 & \textbf{终态效应}：产生伴峰、多重分裂、弛豫位移 \\
电子向表面输运并逸出 & 非弹性散射背景、等离子激元损失峰、逃逸深度限制 \\
\bottomrule
\end{tabularx}}
\end{center}

\subsubsection{初态效应与化学位移的经验规律}

初态即光电发射前原子的基态。当原子与其他原子成键，初态能量改变，结合能随之改变，
此变化即化学位移。

\begin{keycond}
\textbf{一级近似}：同一元素的所有芯能级产生\textbf{相同}的化学位移。
例如 Si 与 O 成键生成 Si—O 后，Si 2p 与 Si 2s 峰相对单质 Si 的位移相似。
这条近似的实用价值在于：它让你可以用一个「内标」核对位移的合理性——
若同一元素的 2p 与 2s 位移方向相反或量级差很多，应先怀疑荷电、拟合或标定问题，
而不是解释成两种化学态。
\src{S2 §9.2.5，书内 p.616 / PDF p.107}
\end{keycond}

对大多数样品，氧化程度增加使结合能升高，用初态效应解释化学位移是合理的。
但下列\textbf{经验规律}都不是定量定律，只能作为判据的辅助
（原书列为四条，见 \src{S2 §9.2.5，书内 p.616--617 / PDF p.107--108}）：

\begin{enumerate}
\item 位移随元素化合价升高而增加——\textbf{同周期主族元素符合，过渡金属不符合}
（如 Cu、Ag），这正是后文必须借助伴峰与俄歇峰的原因。
\item 位移量与同该原子相结合的原子的电负性之和有一定线性关系。
\item 对某些化合物，位移与 NMR、Mössbauer 测得的特征位移量有一定线性关系。
\item 位移与宏观热化学参数有一定联系（原书图 9-59 给出 C 1s 实测值与热力学估算能量的关系）。
\end{enumerate}

\begin{pitfall}
把化学位移直接等同于价态，是本方法最常见的过度断言。三个必须排除的替代解释：
\textbf{（1）荷电效应}——绝缘样品表面荷电会使全部峰同向平移，与真实化学位移无法仅凭
峰位区分，必须先做标定（常用污染碳 C 1s 或 Au 4f$_{7/2}$）；
\textbf{（2）终态效应}——见下节，伴峰与弛豫位移会叠加在主峰上，使峰位与峰形同时改变；
\textbf{（3）配位数与晶体结构变化}——即使价态不变，配位环境改变也会带来零点几 eV 的位移，
与低量级价态变化难以分辨。结论至少需要参考态实测谱与另一独立手段（如俄歇参数）支持。
\end{pitfall}

\subsection{终态效应及伴峰}
\label{sec:ch09-satellites}

\begin{probout}
电离过程产生的各种激发使体系处于\textbf{不同的终态}，不同终态对应不同的结合能，
其对谱的影响统称\textbf{终态效应}。它在谱图上的表现是：除了正常光电子主峰之外，
还出现若干\textbf{伴峰}。伴峰既是干扰（分峰时容易被误当成独立物种），
也是信息（其结构与电子态、自旋态密切相关，可用于区分主峰位移区分不了的化学态）。
\end{probout}

多电子激发的总概率可以很高：每吸收一个光子出现多电子激发过程的总概率
\textbf{可高达 $20\%$}，其中两电子过程的概率大致是三电子过程的 $10$ 倍
\src{S2 §9.2.6，书内 p.617 / PDF p.108}。因此伴峰不是可以忽略的弱效应。

图~\ref{fig:satellites} 把四类伴峰在一张谱上的相对位置画在一起——这是本节最该记住的读图框架：
\textbf{结合能方向是向右增大的}，所以所有「丢失了能量」的过程都出现在主峰右侧。

\begin{figure}[htbp]
\centering
\begin{tikzpicture}
\begin{axis}[
  width=\onecolwidth, height=52mm,
  xlabel={结合能 $E_{\mathrm{B}}$ / \unit{eV}}, ylabel={强度 (a.u.)},
  xmin=278, xmax=326, ymin=0, ymax=1.15,
  xtick={280,290,300,310,320}, ytick=\empty,
  legend pos=north east, legend cell align=left,
  declare function={
    pk(\x,\c,\w,\h)=\h/(1+((\x-\c)/\w)^2);
  }]
% 主峰 + 震激 + 两级等离子激元损失 + 震离台阶（模型计算，非实测）
\addplot[domain=278:326,samples=900,mainline,black]
  {pk(x,285,0.75,1.00)+pk(x,291.5,1.30,0.16)+pk(x,300,1.60,0.26)
   +pk(x,315,2.00,0.12)+0.10*0.5*(1+(x-288)/sqrt((x-288)^2+9))};
% 参考基线
\addplot[domain=278:326,samples=2,thinline,black,dashed]{0};
\legend{模型计算谱,基线}
% 标注
\node[anchor=south,font=\fontsize{8}{9.6}\selectfont] at (axis cs:285,1.02) {主峰};
\draw[thinline,accent] (axis cs:291.5,0.30) -- (axis cs:291.5,0.52);
\node[anchor=west,font=\fontsize{8}{9.6}\selectfont,align=left]
  at (axis cs:292.5,0.52) {震激（卫星）峰\\高结合能侧数 eV};
\draw[thinline,accent] (axis cs:300,0.26) -- (axis cs:300,0.70);
\node[anchor=west,font=\fontsize{8}{9.6}\selectfont,align=left]
  at (axis cs:301,0.70) {等离子激元损失峰\\约 $+5$--$20$ eV，可有多级};
\draw[thinline,accent] (axis cs:288,0.12) -- (axis cs:288,0.30);
\node[anchor=west,font=\fontsize{8}{9.6}\selectfont,align=left]
  at (axis cs:280,0.34) {震离台阶\\连续本底抬高};
\end{axis}
\end{tikzpicture}
\caption{四类伴峰在 XPS 谱上的相对位置。读者应看\textbf{方向}：结合能向右增大，
所以震激峰、等离子激元损失峰都出现在主峰\textbf{右侧}（它们都使出射光电子动能下降）；
震离表现为本底台阶而不是峰；等离子激元损失可以出现多级（间距等距）。
本图为按各过程典型能量间距计算的示意谱，不代表任何具体样品或原书实测图。
\modelmark\figsrc{各类伴峰的相对位置与判据据 S2 §9.2.6（书内 p.617—623 / PDF p.108—114）整理绘制。}}
\label{fig:satellites}
\end{figure}

\subsubsection{震激峰与震离峰}

\textbf{震激（shake up）}是光激发与光电离同时发生的现象：内层电子被光电发射后，
终态弛豫（退激发）导致另一壳层电子被激发到更高能层。它属\textbf{单极激发}——
只有主量子数改变，角量子数与自旋量子数均不变。震激消耗能量，使出射光电子动能下降，
因此震激峰出现在\textbf{主峰高结合能侧、离主线约几个 eV}处。

\textbf{震离（shake off）}是价壳层电子跃迁到非束缚的连续态成为自由电子。它与震激同属
单极激发、同样消耗能量，但在谱图上表现为\textbf{从某个能量开始的连续本底台阶}，
而不是一个峰。震离信号极弱、常被淹没在背底中，实际测试中常被忽略。

\begin{keycond}
震激峰的价值在于\textbf{区分单靠主峰位移区分不了的化学态}。最常被引用的例子是
Cu/CuO/\ce{Cu2O} 体系：用通常的结合能位移区分三者很困难，但它们的 Cu 2p 震激峰
差别明显——\textbf{Cu 和 \ce{Cu2O} 没有明显的 2p$_{3/2}$ 震激峰，而 CuO 有明显的震激峰}。
另一条规律：震激现象在\textbf{顺磁}化合物中发生得很普遍；稀土元素的震激峰强度甚至可以
与主峰相比拟（四价 Ce 的 3d 谱中 3d$_{5/2}$ 与 3d$_{3/2}$ 各有两个震激峰，
且\textbf{震激峰强度超过光电子主峰}）。
\src{S2 §9.2.6，书内 p.618--619 / PDF p.109--110}
\end{keycond}

\subsubsection{X 射线激发的俄歇峰与等离子激元损失峰}

用 \qtyrange{100}{3000}{eV} 的 X 射线源时，绝大多数元素除出射光电子外还出射俄歇电子，
因此 XPS 谱中常伴有俄歇跃迁谱线。俄歇峰的动能同样与化学环境密切相关，
可用其化学位移与峰形变化研究化学态。\textbf{由于俄歇跃迁涉及三电子过程，
其化学位移量往往比光电子峰大得多}——当光电子峰位移不明显时，俄歇峰的位移与峰形
可以提供补充信息。典型例子是 Zn 与 ZnO：Zn 2p 谱较难区分，而 X 射线激发的
Zn LMM 俄歇谱线在峰位与峰形上差别都较大。

\textbf{等离子激元能量损失峰}来自光电子逃逸过程中的非弹性散射。当光电子能量在
\qtyrange{100}{1500}{eV} 时，主要的非弹性散射方式是激发固体中自由电子的集体振荡
（等离子激元）。与一般非弹性散射不同，等离子激元的能量损失是\textbf{特征的、量子化的}。
金属样品通常在光电子峰低动能端 \qtyrange{5}{20}{eV} 处出现主损失峰，并伴随一系列
等间距次级峰；非导体则多表现为一系列长拖尾峰。

\begin{table}[htbp]
\centering
\caption{四类伴峰的来源、位置与用途\src{据 S2 §9.2.6，书内 p.620--622 / PDF p.111--113 归纳。}}
\label{tab:satellite}
{\srcsize
\begin{tabularx}{\columnwidth}{@{}lXX@{}}
\toprule
类型 & 来源与位置 & 读图要点 \\
\midrule
震激（卫星）峰 & 弛豫使另一壳层电子被激发；主峰高结合能侧数 eV & 单极激发；顺磁体系显著；可区分 Cu/\ce{Cu2O}/CuO 类化学态 \\
震离台阶 & 价电子跃迁到连续态；从某能量起本底台阶 & 不是峰而是台阶；信号极弱，常被忽略 \\
俄歇峰 & X 射线激发的三电子过程 & 位移量常大于光电子峰；Zn/ZnO 类体系的关键补充 \\
等离子激元损失 & 逃逸时激发自由电子集体振荡；主峰低动能端 $5$--$20$ eV，可多级 & 金属为等间距尖峰，非导体为长拖尾 \\
\bottomrule
\end{tabularx}}
\end{table}
\emph{核对说明}：已对照原页面图像（书内 p.620--622 / PDF p.111--113）。原书给出：
震激（卫星）峰与震离台阶的区别——震离是价壳层电子跃迁到非束缚的连续状态，从某个能量开始出现连续的较高本底台阶（图
9-64、图 9-65 为示意）；X 射线激发的俄歇跃迁谱线——在 100$\sim$3000 eV 的 X
射线源下绝大部分元素除出射光电子外还出射俄歇电子，俄歇峰的动能也与元素所处的化学环境有密切关系，
其化学位移量往往比光电子峰的更大（图 9-66 为 Zn 与 ZnO 的 Zn 2p 谱及 X 射线激发的 Zn
LMM 俄歇谱，峰位与峰形差别较大）；等离子激元能量损失峰——金属通常在光电子峰的低动能端 5$\sim$20
eV 处观察到主要损失峰，还有一系列次级峰（图 9-67 为金属铝的 2s 光电子峰与一系列的能量损失峰），
而对于非导体通常是一系列长拖尾峰。表中四行均可由这三页支持。

\begin{pitfall}
\textbf{（1）把伴峰当成独立物种。}震激峰与等离子激元损失峰都是同一元素主峰的伴随结构；
若在分峰时把它们当作额外的化学态，会凭空多出并不存在的组分。
判断方法：伴峰与主峰的能量间距是否恒定、是否随激发源改变、是否与主峰强度成比例。
\textbf{（2）分峰时把伴峰面积并入主峰。}反向的错误同样致命：定量时若把震激峰计入主峰面积，
该元素含量会被系统性高估。正确做法是明确交代哪些峰被计入、哪些被排除。
\textbf{（3）忽略等离子激元损失对厚样品的贡献。}损失峰强度随样品厚度增加，
在金属厚膜上可以强到干扰邻近元素的定量。
\end{pitfall}

\begin{crosscheck}
本节的判据都依赖\textbf{参考态实测谱}而非纯理论。若某化学态的归属对结论关键，
应同时给出：（1）参考化合物的实测谱；（2）改变化学态的对照实验（如原位还原/氧化后复测）；
（3）独立手段的交叉验证——表面元素与化学态的整体框架见分片 a，
\textbf{俄歇电子能谱（AES）}的独立介绍见 \chapref{ch:09} 后续分片，
结构层面的对照应用 XRD（\chapref{ch:04}）与化学吸附（\chapref{ch:05}）。
\end{crosscheck}

\srcnote{本分片来源：S2（下册）§9.2.5--§9.2.6，书内 p.615--623 / PDF p.106--114。
式~\eqref{eq:eb-def} 对应原书对 $E_{\mathrm{B}}$ 的定义（OCR 作
「Eb = Ef(n1)Ei(n)」，括号与减号均丢失，按幂次与上下文重建为
$E_f(n-1)-E_i(n)$）。
原书本节所配的实测谱图（图 9-58、9-60—9-68）与震激峰实例均为实验谱，
本书\textbf{未裁切复制}，改以模型计算的示意谱（图~\ref{fig:satellites}）说明相对位置关系；
需要看真实谱形与时例的读者请按上述页码回查原书。
9.2.7—9.2.10（其他初级结构谱线、定量分析、数据后处理、常见问题）由分片 d 续写。}
